% Auto-converted from Superweights.docx
% Compile with XeLaTeX.
\documentclass[11pt]{article}
\usepackage[margin=1in]{geometry}
\usepackage{graphicx}
\usepackage{hyperref}
\usepackage[normalem]{ulem}
\newcommand{\ul}[1]{\uline{#1}}
\usepackage{amsmath,amssymb}
\usepackage{fontspec}
\setmainfont{DejaVu Serif}
\usepackage{xcolor}
\usepackage{enumitem}
\usepackage{longtable,booktabs,array}
\usepackage{caption}
\usepackage{float}
\graphicspath{{media/}}
\hypersetup{colorlinks=true,linkcolor=blue,urlcolor=blue,citecolor=blue}
\setlength{\parindent}{0pt}
\setlength{\parskip}{6pt}
\begin{document}
\textbf{Super-Weight Phenotypes Across Genomic Language Model
Architectures}

\textsuperscript{1} Division of Pharmacology and Toxicology, College of
Pharmacy, The University of Texas at Austin, Dell Paediatric Research
Institute, Austin, TX, USA.

*correspondence to:
\href{mailto:ilias@austin.utexas.edu}{\ul{ilias@austin.utexas.edu}}

\textbf{Abstract}

\textbf{Introduction}

Genomic language models use large-scale self-supervised learning to
decode the complex, overlapping vocabularies of DNA sequence and predict
regulatory, functional, and disease-relevant genomic effects
\href{https://paperpile.com/c/FLFbSd/Fbg9+hmJS+Tixt}{\textsuperscript{1--3}}.
Genomic language models now span diverse architectural families,
including autoregressive transformer decoders such as GENERator
\href{https://paperpile.com/c/FLFbSd/wsei}{\textsuperscript{4}} and NTv3
\href{https://paperpile.com/c/FLFbSd/zg9K}{\textsuperscript{5}},
bidirectional encoder models such as DNABERT-2
\href{https://paperpile.com/c/FLFbSd/5VVM}{\textsuperscript{6}},
state-space and convolutional hybrids such as Evo2
\href{https://paperpile.com/c/FLFbSd/j27Z}{\textsuperscript{7}},
Hyena-attention hybrids such as HybridNA
\href{https://paperpile.com/c/FLFbSd/HqTS}{\textsuperscript{8}}, and
EMA-based architectures such as megaDNA
\href{https://paperpile.com/c/FLFbSd/0vKI}{\textsuperscript{9}}.

Despite their growing predictive success, the internal mechanisms
through which genomic language models encode and understand genomic
information remain largely unresolved. This gap is increasingly
important because genomic language models are being used not only as
predictors, but also as tools for biological discovery, variant
interpretation, model compression, and clinical prioritization. Recent
work in natural language models has identified rare super-weight neurons
whose ablation causes catastrophic model failure despite comprising an
infinitesimal fraction of total parameters
\href{https://paperpile.com/c/FLFbSd/h0C1}{\textsuperscript{10}}.
Whether analogous phenomena exist in genomic language models, and
whether they generalize across fundamentally different sequence-modeling
architectures, remains unknown.

Here, we systematically investigate super-weight phenotypes across six
genomic language models spanning transformer decoders, encoder-only
transformers, state-space models, exponential moving-average
architectures, and hybrid attention, Hyena systems. We show that
catastrophic single-row dependencies emerge specifically in
transformer-decoder genomic models, where individual down-projection
rows controlling \textless0.001\% of parameters can increase perplexity
by more than 25,000\% upon ablation, whereas alternative architectures
distribute representational burden more diffusely and exhibit no
comparable fragility.

\textbf{Results}

\textbf{Super-weight behavior is specific to transformer-decoder genomic
LMs}

We applied the iterative activation-sweep algorithm of Yu et al. (2024)
\href{https://paperpile.com/c/FLFbSd/h0C1}{\textsuperscript{10}} to six
publicly available genomic language models spanning four architectures:
GENERator eukaryote 3B and prokaryote 3B (Llama-based autoregressive
decoders)
\href{https://paperpile.com/c/FLFbSd/wsei}{\textsuperscript{4}},
DNABERT-2 (BERT encoder)
\href{https://paperpile.com/c/FLFbSd/5VVM}{\textsuperscript{6}},
Nucleotide Transformer v3 (NTv3; Mistral-based decoder)
\href{https://paperpile.com/c/FLFbSd/zg9K}{\textsuperscript{5}}, Evo2 7B
(StripedHyena2 SSM/CNN hybrid)
\href{https://paperpile.com/c/FLFbSd/j27Z}{\textsuperscript{7}},
HybridNA (Hyena-attention hybrid)
\href{https://paperpile.com/c/FLFbSd/HqTS}{\textsuperscript{8}}, and
MegaDNA (EMA-based Mega architecture)
\href{https://paperpile.com/c/FLFbSd/0vKI}{\textsuperscript{9}}. For
each model we recorded the per-layer maximum hidden-state activation
across a representative probe set, identified any row in the down\_proj
weight matrix whose zeroing caused the largest activation drop, and then
evaluated the functional impact of zeroing that row by measuring
perplexity on held-out sequences (\textbf{Figure 2a}); for DNABERT-2 we
additionally report accuracy on three Genomic Understanding Evaluation
(GUE) benchmark tasks.

In GENERator-Eukaryote 3B, the sweep converged on a single dominant
super-row at layer four, row 2,371 (\textbf{Figure 1a})(peak activation:
out\_max = 375,361; in\_max = 67,551). Although this row contains only
8,448 parameters, approximately 0.0003\% of the model's \textasciitilde3
billion weights, ablating it increased perplexity from 3.40 to 785.5
(+23,026\%) (\textbf{Figure 2b}). By contrast, ten randomly selected
rows of identical size produced no measurable effect (ΔPPL \textless{}
0.01\%), demonstrating that the collapse is highly specific to the
identified row rather than a generic consequence of parameter removal. A
secondary super-row at the same layer (row 1,522; (\textbf{Figure
1a});out\_max = 143,880) produced an additional but smaller increase
when ablated sequentially, indicating that the primary row accounts for
most of the functional dependency. A similarly extreme phenotype emerged
in GENERator-Prokaryote 3B, where the activation sweep identified a
super-row at layer two, row 1,927 (\textbf{Figure 1b})(out\_max =
506,014). Zeroing this single row increased perplexity from 2,324 to
606,073 (+25,975\%), whereas matched random-row controls again remained
within 0.03\% of baseline performance (\textbf{Figure 2b}). Relative to
the eukaryote model, the earlier layer position and larger activation
magnitude suggest a more globally integrated residual-stream component
(see Mechanistic Interpretability section).

\includegraphics[width=6.26772in,height=8.11111in]{media/image5.png}

\textbf{Figure 1. Super-weight behaviour across genomic language
models.}

\textbf{(A)} Peak activation magnitude (out\_max) of the dominant weight
row plotted against its normalised layer position for seven genomic
language models. Red: models in which zeroing the row raises perplexity
by \textgreater5\% (SW-positive); blue: SW-negative models. \textbf{(B)}
Perplexity change (ΔPPL\%) after ablating the super-weight row (bars)
versus an equal-norm random row (dots); grey band marks the
random-control envelope. \textbf{(C--D)} Location of the super-weight
within the GENERator-EUK-3B and GENERator-PROK-3B transformer stacks.
The highlighted layer (blue) contains the super-weight; orange stripes
in the down\_proj weight matrix indicate the affected row(s).
\textbf{(E--F)} Activation lifecycle of the super-weight channel across
all 30 layers. Top: post-residual hidden-state magnitude (mean ± 1 s.d.
over 20 sequences); bottom: per-block contribution to that channel.
Vertical dotted line marks the super-weight layer. Grey band in the top
panels shows the range spanned by the top-3 globally active channels as
a reference.

DNABERT-2 showed a qualitatively different pattern. The activation sweep
identified ten potential super-rows concentrated around row 603 at
multiple encoder layers (layers 3, 5, 6, 7, 9; out\_max ranging from 137
to 945). In isolation, zeroing these rows produced only a modest
perplexity increase of +1.5\% (baseline 4.486 versus ablated 4.554)
(\textbf{Figure 2b}), attributable to DNABERT-2\textquotesingle s
masked-language-model pretraining objective, which averages loss over
all tokens and partially masks single-row effects. To assess whether SW
rows encode functionally relevant representations, we ablated all 10 SW
rows of DNABERT-2 and evaluated downstream performance on three GUE
benchmarks across three random seeds. Splice-site recognition showed a
robust and consistent drop of −25.5 ± 0.7 percentage points (p = 0.0004,
one-sample t-test vs. zero), confirming that SW rows are required for
this task. Promoter recognition and histone modification prediction
(H3K4me3) showed mean drops of −12.6 ± 16.7 pp (p = 0.40) and −4.5 ± 6.2
pp (p = 0.41) respectively; these effects are not statistically
significant at n = 3 seeds and should be interpreted as exploratory. In
all three tasks, ablating an equal number of randomly selected rows
produced changes of \textless{} 0.03 pp, confirming that the effect is
specific to SW rows.

By contrast, NTv3, Evo2, HybridNA, and MegaDNA showed no convincing
evidence of super-weight behavior. NTv3 represented a borderline case:
its top activation row (L11, r1472) produced a maximum ΔPPL of +4.8\% on
one probe sequence, but the effect was inconsistent across inputs
(−2.1\% to +4.8\%) and indistinguishable from matched random-row
controls (\textbf{Figure 2a-b}). Evo2, HybridNA, and MegaDNA exhibited
essentially no functional sensitivity to row ablation (ΔPPL = +0.0007\%,
−1.3\%, and +0.3\%, respectively), despite Evo2 containing activation
outliers comparable in magnitude to those observed in GENERator
(out\_max = 552,960 at layer 29) (\textbf{Figure 2a-b}). These findings
indicate that extreme activations alone are insufficient to generate
super-weight phenotypes. Instead, the phenotype may be linked to
architectural differences in information routing: transformer decoders
concentrate information through residual-stream bottlenecks and
down-projection pathways, whereas SSM-, EMA-, and Hyena-based
architectures distribute representations more diffusely across recurrent
or convolutional filters. Collectively, these results suggest that
super-weight behavior can be a specific property of transformer-decoder
geometry rather than a universal consequence of scale or genomic
sequence modeling.

\textbf{Super-weights encode sequence composition, not regulatory
context}\\
To determine whether the identified super-rows are causally necessary
for model function rather than merely exhibiting large activations, we
zeroed each row during inference and measured the resulting perplexity
increase on 90 held-out genomic sequences (30 each from promoters,
enhancers, and intergenic regions for the eukaryote model; promoters,
terminators, and random regions for the prokaryote model). In the
eukaryote model, ablating row 2,371 at layer four increased perplexity
by a mean of +2.92 log-PPL units (promoter: +2.68, enhancer: +2.79,
random: +3.28). In the prokaryote model, ablating row 1,927 at layer two
caused a substantially larger increase of +9.47 log-PPL units (promoter:
+9.15, terminator: +9.34, random: +9.93). In both cases, a randomly
chosen control row produced a mean shift of \textless0.0001 log-PPL
units. The context-uniform degradation across all three genomic classes
establishes that the super-rows are not specialised for any regulatory
element, meaning they are general-purpose components of the residual
stream.

\includegraphics[width=6.26772in,height=8.11111in]{media/image4.png}

\textbf{Figure 2. Super-weight functional specificity revealed by causal
tracing, sequence composition, and task-selective activation.}

\textbf{(A--B)} Causal tracing results for GENERator-EUK-3B \textbf{(A)}
and GENERator-PROK-3B \textbf{(B)}. Bars show the mean change in
per-token perplexity (ΔPPL) when the super-weight (SW) row or a matched
random row is ablated, across promoter, enhancer, and random genomic
contexts \textbf{(A)} or promoter, terminator, and random genomic
contexts \textbf{(B)}. Error bars indicate ±1 SD over ablation
replicates. \textbf{(C--D)} Mean absolute super-weight activation
grouped by hexamer sequence category for GENERator-EUK-3B \textbf{(C)}
and GENERator-PROK-3B \textbf{(D)}. In the eukaryotic model
\textbf{(C)}, hexamers are grouped by functional class: splice-donor
(GT-starting), splice-acceptor (AG-ending), AT-rich (≥5 AT bases), and
other. The same sequence-based grouping is applied to the prokaryotic
model \textbf{(D)}, where categories reflect sequence composition rather
than splicing function. Error bars show ±1 SEM; n indicates the number
of hexamers per group. \textbf{(E--F)} Motif enrichment analysis of the
top 500 hexamers ranked by KL divergence between clean and ablated model
outputs. Bars show the odds ratio for enrichment of each motif class
among high-KL hexamers; red bars indicate statistically significant
enrichment after Benjamini--Hochberg correction (q \textless{} 0.05),
grey bars are non-significant (legend, panel F). \textbf{(G--H)}
Super-weight task-activation difference heatmaps for GENERator-EUK-3B
\textbf{(G)} and GENERator-PROK-3B \textbf{(H)}, showing the mean
super-weight activation on functional (positive) sequences minus
non-functional (negative) sequences for each task--super-weight-row
combination. Cell annotations report the mean difference ± pooled SE.
Diverging colour scale is symmetric around zero; positive values (red)
indicate higher activation on functional sequences.

\textbf{Sequence context-independence}\\
To test whether the super-row activation is sensitive to sequence
context above the single k-mer level, we applied four increasingly
strict shuffle controls (mononucleotide, dinucleotide, k-mer block,
trinucleotide) to each of the 90 input sequences and measured the change
in super-row activation. In the eukaryote model, none of the four
shuffles produced a significant activation shift (all p \textgreater{}
0.31, n = 90); the strictest control (trinucleotide-preserving shuffle)
yielded p = 0.84. The SW activation is therefore insensitive to sequence
context at any level above individual hexamer identity. In the
prokaryote model the picture is reversed: all four shuffle types
significantly perturbed SW activation (mononucleotide p = 0.0017,
dinucleotide p = 0.0092, k-mer block p = 0.0099, trinucleotide p
\textless{} 10⁻⁸), indicating that the prokaryote super-row is sensitive
to sequence composition at multiple levels.

\textbf{K-mer identity and the causal hexamer test\\
}A genome-wide hexamer activation scan (all 4,096 6-mers) confirmed that
the eukaryote super-row is activated by nearly all hexamers (IQR
387K-409K), with a modest positive GC-content correlation (r = +0.19)
and CC/CT-family k-mers at the top of the activation distribution. To
close the causal chain, we measured the KL divergence of the next-token
distribution before and after SW ablation for all 4,096 hexamers: the
Pearson correlation between SW activation and KL divergence was r =
+0.437 (t = 5.06, p = 4.46 × 10⁻⁷), confirming that the hexamers that
most strongly activate the SW are precisely those whose predictions are
most disrupted by its removal.

The prokaryote model shows an inverted pattern (GC correlation r =
−0.09; AT-rich hexamers dominate the activation distribution) and a
striking causal paradox: the correlation between SW activation and KL
divergence is r = −0.710 (t = −118, p ≈ 0). The top-quartile AT-rich
activators produce a mean KL of 0.043 after ablation, while the
bottom-quartile low-activation hexamers produce a mean KL of 0.884 which
is more than twenty times larger. The SW is therefore not causally
necessary for its own highest-activating tokens; instead, its removal
catastrophically disrupts predictions for GC-rich hexamers it barely
activates. We interpret this as a suppressive gating mechanism: the
prokaryote SW amplifies AT-rich inputs while simultaneously maintaining
the representational capacity needed to predict GC-rich, compositionally
demanding tokens.

To characterise the sequence features that underlie SW dependency, we
performed Fisher\textquotesingle s exact test with Benjamini--Hochberg
correction comparing the top-500 most SW-dependent k-mers (ranked by KL
divergence) against the remaining 3,596 in both models. In the eukaryote
model, homopolymer-containing k-mers were the only significantly
enriched class among SW-dependent sequences (OR = 1.85, q = 0.022);
canonical regulatory motifs --- splice donor, splice acceptor, TATA-box,
Kozak, CpG-rich, and GC-box --- were all depleted or at odds ratio ≈
1.00 (all q = 1.00). In the prokaryote model, no motif class survived
correction; GT-containing k-mers showed a non-significant trend toward
SW dependency (OR = 1.21, q = 0.33), and Shine--Dalgarno, −10-box, and
−35-box elements showed no enrichment. Together, these enrichment
profiles confirm that SW dependency tracks a sequence-complexity axis
--- structurally repetitive, low-diversity sequences in the eukaryote
model and compositionally mixed, GC-variable sequences in the prokaryote
--- rather than recognisable biological regulatory grammar at the
single-hexamer level.

Taken together, the hexamer causal tests reveal a fundamental divergence
in SW function across kingdoms. In the eukaryote model the SW acts as a
\emph{driver}: high activation predicts high functional load (r =
+0.437), and the sequences it is most responsible for are low-complexity
homopolymeric runs. In the prokaryote model the SW acts as a
\emph{background normaliser}: high activation on AT-rich k-mers predicts
low dependency (r = −0.710), while GC-rich and compositionally diverse
k-mers --- those the SW barely activates --- are the ones whose
predictions collapse when it is removed. This kingdom-level asymmetry
mirrors the compositional difference between the two training corpora:
eukaryotic genomes are GC-heterogeneous with abundant repetitive
elements that the SW must route consistently, while prokaryotic genomes
present a sharper AT/GC compositional contrast that the SW appears to
resolve through suppressive gating. Critically, neither
model\textquotesingle s SW encodes splice sites, promoter elements, or
any other annotated regulatory signal at the k-mer level; the causal
role of the SW is real and statistically robust, but its substrate is
nucleotide composition and sequence complexity rather than biological
function.

\textbf{Sequence composition explains partial but not complete
variance}\\
Multivariate regression of ΔPPL on five sequence composition features
(GC fraction, CpG density, k-mer entropy, repeat fraction, complexity)
explained R² = 0.158 of variance in eukaryote ablation impact, with
sequence complexity (β = −1.48) and k-mer entropy (β = +1.05) as the
dominant predictors. The prokaryote model showed higher explained
variance (R² = 0.347), with GC fraction (β = −0.306) and k-mer entropy
(β = −0.307) both negatively predicting impact consistent with the
observation that AT-rich, low-complexity sequences are most disrupted
when the prokaryote SW is removed.

\textbf{Super-weight activations are not kingdom-specific}\\
Finally, we tested whether each model\textquotesingle s super-row fires
differently on sequences from the opposing biological domain. The
eukaryote super-row (layer 4, row 2371) showed a mean activation of
6,582 on hg38 sequences and 6,610 on E. coli sequences (fold-change =
1.004, Mann-Whitney p = 0.538). The prokaryote super-row (layer 2, row
1927) showed 10,638 on E. coli and 10,710 on hg38 (fold-change = 1.007,
p = 0.644). Neither super-row distinguishes between eukaryotic and
prokaryotic sequence input. Combined with the k-mer identity and
composition results, this establishes that both super-rows encode
low-level statistical features of nucleotide composition present across
all genomes rather than kingdom- or element-specific biological signals.

\includegraphics[width=6.26772in,height=3.13889in]{media/image1.png}

\textbf{Figure 5:}

Violin/swarm plots comparing SW activation distributions when each model
is fed in-domain versus out-of-domain sequences (eukaryote model on E.
coli; prokaryote model on hg38). Both fold-changes are \textless0.7\%
and Mann-Whitney p \textgreater{} 0.5, showing that super-weights are
not kingdom-specific detectors.

To further test whether SW activation encodes functional sequence
identity, we measured SW row activation magnitude on positive-class
versus negative-class sequences across two complementary benchmark
suites. For GENERator EUK, we evaluated eleven GUE tasks spanning
promoter prediction, splice-site recognition, histone modification,
transcription factor binding, and viral sequence classification. For
GENERator PROK, we evaluated five prokaryotic functional categories
derived from the E. coli K-12 MG1655 reference genome: promoters,
terminators, tRNA genes, rRNA genes, and coding sequences (CDS). In both
models, SW activation was indistinguishable between positive-class and
negative-class sequences within every task (fold-change range
0.984--1.013 for EUK rows 2371 and 1522; range 0.997--1.004 for PROK row
1927; Mann-Whitney p \textgreater{} 0.2 in all cases). While absolute
activation magnitudes varied across task categories --- EUK promoter
tasks produced substantially higher values (38K) than histone or viral
tasks (2.5--5K), and PROK activations were uniform across all five
categories (\textasciitilde8.2K) --- this variation reflects differences
in sequence length and tokenization depth rather than biological
content. No task in either model produced a consistent positive/negative
class distinction (Figure X). These results confirm that SW rows
implement a sequence-length-scaled global amplification operation rather
than a class-discriminative feature detector, and that their causal
necessity for downstream tasks arises through a mechanism other than
differential firing on functional versus non-functional sequences.

\textbf{Super-weights arise from a quadratic amplifier in the step-up
FFN block}

The preceding sections establish that SW rows are detectable,
functionally critical, and sequence-composition-sensitive. Here we ask a
more fundamental question: why do these specific weights acquire extreme
activation magnitudes? We address this mechanistically, following the
framework of Sun et al. (2026)
\href{https://paperpile.com/c/FLFbSd/XSIF}{\textsuperscript{11}}, who
showed that massive activations in NLP language models arise from a
"directional quadratic amplifier" structure within a single early
feed-forward block. We show that the same mechanism operates in genomic
LMs.

Activation lifecycle. For each model we registered hooks on every
Transformer layer and tracked the maximum activation magnitude in the SW
channel across all 30 layers, averaged over 20 genomic sequences. In
both GENERator EUK and PROK, the post-residual SW channel magnitude
follows a characteristic rise--plateau--fall trajectory (Figure X).
Prior to the step-up layer, the SW channel magnitude is
indistinguishable from the global background (EUK: \textasciitilde300;
PROK: \textasciitilde850). At the step-up layer --- layer 4 in EUK and
layer 2 in PROK --- the channel magnitude jumps by three orders of
magnitude in a single forward pass (EUK: 375,000; PROK: 509,000). The
block-level delta at the step-up layer accounts for \textgreater99.9\%
of the total SW channel magnitude, confirming that the injection is
abrupt rather than cumulative. Following injection, the massive
activation persists essentially unchanged through all subsequent
intermediate layers via the residual stream; intermediate-layer block
contributions to the SW channel are two to three orders of magnitude
smaller than the activation itself. A partial step-down signal is
visible in the final one to three layers (EUK: delta = 5,081 at layer
29; PROK: delta = 257,000 at layers 27--29), consistent with late-layer
neutralisation observed in NLP models (Sun et al., 2026). This
rise--plateau--fall pattern directly explains the fragility of SW rows
to ablation: the activation functions as a persistent global signal that
downstream layers integrate across the full depth of the network.

Quadratic amplifier fingerprint. To identify the structural origin of
the step-up injection, we computed the Frobenius norm of the effective
quadratic form ‖U\_k‖\_F for every output coordinate k of the FFN at the
step-up layer, where U\_k summarises the combined amplification capacity
of the gate, up, and down projection weights feeding into output row k
(Sun et al., 2026; Yu et al., 2024). Under the near-identity gating
regime that holds for spike tokens, ‖U\_k‖\_F is proportional to the
maximum activation that output coordinate k can produce. SW rows ranked
first among all 3,072 output coordinates in both models (Figure Y): EUK
row 2371 reached ‖U\_k‖\_F = 215.86 versus a median of 17.63 (12× above
median; rank 1/3,072), and EUK row 1522 ranked second (4.6× above
median). PROK row 1927 showed an even more extreme separation, with
‖U\_k‖\_F = 2,648 versus a median of 148 (18× above median; rank
1/3,072). Crucially, the SW rows were identified independently in
earlier experiments by their activation magnitude during inference ---
the fact that they simultaneously top the weight-space mechanistic
ranking confirms that detection and mechanism are measuring the same
underlying structure.

\includegraphics[width=6.26772in,height=5.44444in]{media/image2.png}

\textbf{Figure 8: Frobenius norm ‖U\_k‖\_F audit (change/remove the left
plots).} Quadratic amplifier strength ‖U\_k‖\_F for every output
dimension k of the FFN at the step-up layer in GENERator EUK 3B (top)
and PROK 3B (bottom). SW rows (red dashed lines) rank first among all
3,072 output coordinates in both models, with ‖U\_k‖\_F 12× (EUK) and
18× (PROK) above the genome-wide median. SW rows were identified
independently from activation magnitudes during inference; their top
ranking in the weight-space audit confirms that both measures reflect
the same high-gain amplifier structure.

Together these results establish that SW rows in genomic LMs are not
arbitrary outliers but the predictable output of a high-gain FFN channel
whose intermediate dimensions exhibit anomalously collinear gate and up
projections paired with large down-projection entries --- the same
quadratic amplifier architecture identified in NLP LLMs by Sun et al.
(2026), here documented for the first time in a genomic sequence
modelling context.

\textbf{Implications for Model Compression}

\textbf{Pruning: SW-proximal rows are compression-tolerant (DNABERT-2)}

We performed a progressive row-pruning sweep on DNABERT-2 fine-tuned on
the core-promoter prediction task (prom\_core\_notata; baseline accuracy
83.46\%, MCC 0.670), zeroing increasing fractions of down\_proj rows
(0.5\%--30\%) selected by five criteria: L1-low (smallest weight norm,
expected most dispensable), L1-high (largest weight norm),
proximity-near (closest rows in layer/row space to the super-weight
coordinates), proximity-far (most distant), and random.

The key finding is the divergence between prox\_near and prox\_far as
fraction increases. At 20\% of total rows pruned, the near-SW strategy
(prox\_near) retained 98.4\% of baseline accuracy (Δ = −1.60 pp), while
the far-SW strategy (prox\_far) degraded substantially (Δ = −9.39 pp)
--- a 5.9-fold difference driven entirely by proximity to the SW rows.
At 30\% pruning the asymmetry persists: prox\_near loses only −1.65 pp
versus −8.15 pp for prox\_far. Magnitude-based baselines fall between
these extremes: l1\_low reaches −3.93 pp and l1\_high −1.20 pp at 30\%,
confirming that SW proximity is a stronger predictor of row importance
than raw weight magnitude. Rows near the SW can be zeroed with minimal
functional consequence because the SW itself maintains the dominant
representational pathway; rows far from it carry independent signal that
the model cannot recover from ablation.

\includegraphics[width=6.26772in,height=2.44444in]{media/image3.png}

\textbf{Figure 9}: Progressive down-proj row-pruning sweep (5--50\% of
rows) on DNABERT-2 using five selection criteria (SW-fragility, near-SW,
L1-low, random, baseline). Dual y-axes: accuracy and MCC vs fraction
pruned. Shows that SW-adjacent rows are significantly more
compression-tolerant than random rows, which is the core compression
finding for DNABERT-2.

\textbf{Quantization: SW rows are not precision-sensitive (GENERator
INT4)\\
(maybe need to recheck the results)}

To test whether the super-rows identified in GENERator require precision
protection during quantization --- as proposed by Yu et al. (2024) for
NLP LLMs --- we applied round-to-nearest (RTN) INT4 quantization to all
806,396 eligible rows across all seven linear projection types in the
30-layer network, evaluating three conditions: (i) yu\_all: all non-SW
rows quantized, SW rows retained in FP32; (ii) yu\_all\_including\_sw:
all rows quantized including SW rows; (iii) sw\_fragility: only the SW
rows quantized, all others retained.

In the eukaryote model, INT4 quantization of all 806,396 non-SW rows
increased perplexity by +0.117 (baseline 8.421 → 8.537). Adding the 6 SW
rows to the quantization pool (yu\_all\_including\_sw) raised perplexity
by an additional 0.0006 --- a SW marginal cost of +0.0006 PPL.
Quantizing only the SW rows in isolation produced a change of +0.0005.
The prokaryote model showed a larger overall INT4 cost (+0.444 PPL;
baseline 8.697 → 9.141), consistent with the model\textquotesingle s
deeper residual-stream integration demonstrated by causal tracing, but
the SW marginal cost was again +0.0006 (sw\_fragility: +0.0003). In both
models the SW marginal cost is indistinguishable from measurement noise
on a 7,200-token probe.

This is a direct non-replication of the Yu et al. SW-exemption heuristic
in genomic LMs: whereas in NLP transformers retaining SW rows in full
precision is critical to INT4 compression quality, in GENERator the SW
rows are no more precision-sensitive than any other row in the network.
We attribute this to the mechanism established in the interpretability
section: the eukaryote SW row encodes k-mer composition through extreme
input activation magnitude (in\_max = 67,551), not through unusually
large or structured weight values. INT4 quantization of weight matrices
with modest weight magnitude (max weight ≈ 1.85 in down\_proj) is
near-lossless for those rows; it is the ordinary rows, distributed
across all 30 layers and 7 projection types, that accumulate the
dominant quantization error.

\includegraphics[width=6.26772in,height=4.15278in]{media/image6.png}

\textbf{Figure 10. SW rows are not precision-sensitive under INT4
quantization.}

RTN INT4 applied to all 806,396/806,397 linear projection rows
(down/gate/up/q/k/v/o\_proj) across 30 layers. Blue: perplexity
increases when all non-SW rows are quantized. Red: additional perplexity
cost of including the SW rows in the quantization pool. The SW marginal
cost (+0.0006 PPL) is 187× smaller than the overall INT4 cost in the
eukaryote model and 758× smaller in the prokaryote model, demonstrating
that SW rows require no special precision protection --- in contrast to
the SW-exemption heuristic proposed for NLP LLMs.

\textbf{Methods}

Models

We studied three genomic language models. GENERator-EUK-3B
(GenerTeam/GENERator-v2-eukaryote-3b-base) and GENERator-PROK-3B
(GenerTeam/GENERator-v2-prokaryote-3b-base) are causal decoder-only
transformers sharing the same Llama architecture: 30 layers, hidden
dimension

\(d_{model} = 3072\), feed-forward intermediate dimension
\(d_{ffn} = 8448\), and vocabulary size

4,128. DNABERT-2 (zhihan1996/DNABERT-2-117M) is a bidirectional
BERT-style encoder with 12 layers, \(d_{model} = 768\),
\(d_{ffn} = 3072\), 12 attention heads, and vocabulary size

4,096. All models use a byte-pair encoding tokenizer on nucleotide
sequences and were loaded in full precision from HuggingFace Hub with
transformers.

Superweight identification

We define the down-projection matrix
\(W_{down}\epsilon\ R^{d_{model} \times d_{ffn}}\) of each MLP block as
the linear map from the gated intermediate space back to the residual
stream. A superweight row is a row

\(k\) of \(W_{down}\) whose mean absolute value greatly exceeds the
layer-wide average. For GENERator models we inspect
model.layers.\{i\}.mlp.down\_proj; for DNABERT-2 we inspect
model.encoder.layer.\{i\}.intermediate.dense (the mlp.wo equivalent).

To quantify the structural importance of row \(k\), we compute a rank-1
Frobenius bound that captures the amplification potential through the
full MLP block:

\(\left| U_{k} \right|_{F} =\)

where \(W_{gate}\lbrack i,:\rbrack\) and \(W_{up}\lbrack i,:\rbrack\)
are the \(i\)-th rows of the gate and up-projection matrices. This bound
is tight when the outer products of the gate and up rows are aligned,
and is a conservative approximation otherwise.

Causal ablation experiment

To test whether identified superweight rows are functionally
indispensable, we perform a causal ablation (zeroing) experiment. A
PyTorch forward hook is registered on the down\_proj output and zeros
the superweight row(s) in-place on every forward pass; the hook is
removed in a final block to guarantee clean restoration.

Model performance is measured as mean per-token cross-entropy loss in
nats, returned directly by the HuggingFace CausalLM forward pass with
labels=input\_ids (i.e., output.loss.item()). We refer to this quantity
as the loss throughout; it is the negative log-likelihood per token and
is not exponentiated to perplexity. We sample 30 non-overlapping
3,072-bp windows from the E. coli K-12 chromosome (GENERator-PROK) or
from eukaryotic reference sequences (GENERator-EUK) and report the mean
loss over those windows.

Shuffle controls

To establish whether SW activations respond to sequence content rather
than position, we implemented four shuffle controls:

\begin{enumerate}
\def\labelenumi{\arabic{enumi})}
\item
  Mono-nucleotide shuffle (mono): random permutation of all nucleotides,
  destroying all
\end{enumerate}

\(k > 1\) context.

\begin{enumerate}
\def\labelenumi{\arabic{enumi})}
\setcounter{enumi}{1}
\item
  Dinucleotide-preserving shuffle (dinuc): Altschul--Erickson algorithm,
  preserving exact dinucleotide frequencies.
\item
  Trinucleotide-preserving shuffle (trinuc): generalization of
  Altschul--Erickson to trinucleotides, preserving 3-mer counts.
\item
  6-mer block shuffle (kmer\_block): shuffles at the hexamer token
  level, preserving hexamer identity while destroying positional
  context.
\end{enumerate}

Each control is applied to the same sequences used for the causal
ablation.

Hexamer activation scan

The genome-wide hexamer scan is implemented in
scripts/interpretability/run\_sw\_hexamer\_causal.py. For each distinct
hexamer token, we construct a minimal 11-token context (CTX\_HALF = 5
flanking tokens) centered on the target position (TARGET\_POS = 5), feed
it through the model with a forward hook on the SW row, and record the
activation. Poly-A (AAAAAA) is assigned token ID 32 (POLY\_A\_TOKEN =
32). KL divergence between ablated and non-ablated next-token
distributions is computed as:\\
\(D_{KL}(p||q) = \sum_{v}^{}p_{v}\log\frac{p_{v}}{q_{v}}\ \ \)with
\(\epsilon = 10^{- 9}\) added to both distributions before
renormalization to ensure numerical stability.

Hexamer motif enrichment

Motif enrichment over the top-ranked hexamers is computed using a
two-sided Fisher exact test with Benjamini--Hochberg correction, testing
the top-500 hexamers by KL divergence against the full hexamer
background. Eukaryotic regulatory motifs tested include splice donors
(hexamers containing GT), splice acceptors (ending in AG with ≥3
pyrimidine bases), TATA boxes (TATA substring), GC boxes (exact GGGCGG),
Kozak consensus (exact GCCACC or GCCGCC), CpG-rich hexamers (≥2 CG
dinucleotides), AT-rich hexamers (≥5 A/T bases), and homopolymers (any
base repeated ≥4 times consecutively). Prokaryotic signal sequences
tested include Shine--Dalgarno motifs (AGGAGG, GAGGAG, or AAGGAG), −10
box variants (TATAAT, TACAAT, TATACT), and −35 box variants (TTGACA,
TTGATA, TTGAAA). All predicates operate on exact substring matching
within the 6-mer string.

Task-specific activation heatmap

BED files were constructed from the NCBI RefSeq GFF annotation for E.
coli K-12 (GCF\_000005845.2), extracting features of type tRNA (86
entries), rRNA (22 entries), and CDS (500 randomly sampled entries).
Promoter and terminator BED files (4,506 entries each) were constructed
from RegulonDB annotation. A random background BED (1,000 entries) was
also included. For each region, the 3,072-bp window centered on the
feature midpoint is tokenized, the SW activation is extracted via
forward hook, and the mean absolute value over the token dimension is
recorded.

Cross-kingdom transfer experiment

For each model (EUK or PROK), 40 sequences of 384 bp are drawn: half
from the in-domain reference (eukaryotic or prokaryotic), half from the
out-of-domain reference. SW activation is measured for each and the
distributions compared.

Activation lifecycle

The layer-by-layer trajectory of superweight channel magnitudes is
traced by taking forward pre-hooks and post-hooks are registered on
every transformer layer to capture the post-residual hidden state
\(H_{i + 1}\) (input to the next layer) after each residual addition.
For each sequence, we extract
\(\max_{t}\left| H_{i}\lbrack:,k\rbrack \right|\)--- the token-wise
maximum absolute value for SW row \(k\ \)at layer \(i\). Results are
averaged over 20 sequences of 504 nt each, with one E. coli probe
sequence and 19 randomly generated sequences.

Sequence composition regression

To assess whether SW activation magnitude is predictable from sequence
composition, we fit an ordinary-least-squares regression of
ablation-induced loss change (\$\textbackslash Delta\$loss) on five
sequence features, computed in
scripts/interpretability/run\_sw\_ablation\_regression.py:

\begin{enumerate}
\def\labelenumi{\arabic{enumi})}
\item
  GC fraction --- fraction of G/C bases
\item
  CpG density --- observed/expected CpG ratio:
  \((N_{CG} \times L)/(N_{C} \times N_{G})\)
\item
  6-mer entropy --- Shannon entropy of the hexamer token frequency
  distribution: \(H = \  - \sum_{k}^{}p_{k}\log p_{k}\)
\item
  Repeat fraction --- fraction of sequence covered by tandem hexamer
  runs
\item
  Linguistic complexity --- number of distinct hexamers observed divided
  by total hexamer positions
\end{enumerate}

All features are computed per-sequence from the same windows used in the
causal ablation.

Pruning sweep (DNABERT-2)

To evaluate the compressibility of DNABERT-2 with respect to
superweights, we sweep pruning fractions in
scripts/compression/run\_compression\_sweep.py. The candidate pool
consists of all rows of all down\_proj weight matrices that are not
identified as superweights. We sweep fractions
\(f\ \epsilon\ \left\{ 0.5,\ 1,\ 2,\ 5,\ 10,\ 15,\ 20,\ 30 \right\}\%\)
of this pool under five ranking criteria:

l1\_low: ascending row L1-norm (smallest-magnitude rows removed first;
"least important" by magnitude)

l1\_high: descending row L1-norm (largest-magnitude rows removed first;
sanity / upper-bound check)

prox\_far: rows furthest from any superweight row in parameter space
(removes the "irrelevant" periphery first)

prox\_near: rows closest to any superweight row (removes the SW
neighbourhood first)

random: uniformly random order, repeated over 10 seeds; mean ± std
reported

Pruning is implemented as zeroing the selected rows. Model accuracy and
MCC on the GUE benchmark (promoter-core task) are evaluated at each
fraction.

INT4 quantization experiment

Whole-model INT4 quantization is implemented in
scripts/compression/run\_whole\_model\_quantization.py using
round-to-nearest (RTN) quantization with per-row scaling:

\(\widehat{r} = \ round(r/{s) \times s,\ s\  = \ \frac{max\ |r|}{7}}_{}\)

where \(s\) is the per-row scale factor and the integer range is
\(\lbrack - 7,7\rbrack\). We quantize all seven projection types:
down\_proj, gate\_proj, up\_proj, q\_proj, k\_proj, v\_proj, o\_proj.
Four conditions are tested:

Baseline: no quantization

yu\_all: all rows quantized except superweight rows

yu\_all\_including\_sw: all rows quantized including superweight rows

sw\_fragility: only superweight rows quantized; all others left in full
precision

Model quality is evaluated on 6 probe sequences × 1,200 nt = 7,200 nt
total.

\textbf{Discussion}

Big-picture points

\begin{quote}
Super-weight--like catastrophic single-row dependencies appear in
GENERator transformer decoders but not in Evo2, HybridNA, MegaDNA, or
DNABERT-2.

In GENERator, ablating a single down-projection row (\textless0.001\% of
parameters) blows up perplexity, while random rows do nothing.

Non-decoder architectures reach similar activation magnitudes but are
functionally robust to single-row ablations.
\end{quote}

Mechanism / what they seem to do

\begin{quote}
Euk SW: broadly GC-biased, composition/complexity channel; strongly
involved in low-complexity, homopolymer-like sequences, not classic
regulatory motifs.

Prok SW: acts more like a suppressive/normalising gate; high activation
on AT-rich hexamers but biggest functional dependency on GC-rich,
mixed-composition k‑mers.

Neither SW looks like a clean detector for promoters, splice sites, or
other annotated elements; it is more about nucleotide composition than
``biology''.
\end{quote}

Architecture / interpretability angle

\begin{quote}
We scan all linear projections (MLP gate/up/down and attention
projections) and only find true super-weights in down-projection rows of
decoder MLP blocks.

Transformer-decoder geometry plus MLP down-projection bottlenecks seem
to be the ingredients for catastrophic SW behavior; SSM/EMA/Hyena models
smear that burden out.

Mechanistic analysis points to a quadratic-amplifier FFN channel that
injects a huge global signal in a single early layer and then rides the
residual stream, explaining both the massive activations and the
fragility to ablation.
\end{quote}

Practical implications / limitations

\begin{quote}
For genomic LMs in applied settings, it may be worth explicitly checking
for SW rows as a robustness/interpretability diagnostic.

INT4 experiments suggest genomic SWs are not uniquely
precision-sensitive, so NLP-style ``protect the SW row'' quantization
heuristics may not directly transfer.

We still only study six models (plus protein models) and a limited panel
of downstream tasks and sequence features, so there is plenty of room
for broader architectures, objectives, and biological readouts.
\end{quote}

\textbf{References}

1. \href{http://paperpile.com/b/FLFbSd/Fbg9}{Avsec, Ž. \emph{et al.}
Effective gene expression prediction from sequence by integrating
long-range interactions. \emph{Nat Methods} \textbf{18}, 1196--1203
(2021).}

2. \href{http://paperpile.com/b/FLFbSd/hmJS}{Dalla-Torre, H. \emph{et
al.} Nucleotide Transformer: building and evaluating robust foundation
models for human genomics. \emph{Nat Methods} \textbf{22}, 287--297
(2025).}

3. \href{http://paperpile.com/b/FLFbSd/Tixt}{Linder, J., Srivastava, D.,
Yuan, H., Agarwal, V. \& Kelley, D. R. Predicting RNA-seq coverage from
DNA sequence as a unifying model of gene regulation. \emph{Nat Genet}
\textbf{57}, 949--961 (2025).}

4. \href{http://paperpile.com/b/FLFbSd/wsei}{Wu, W. \emph{et al.}
GENERator: A Long-Context Generative Genomic Foundation Model. (2025).}

5. \href{http://paperpile.com/b/FLFbSd/zg9K}{Boshar, S. \emph{et al.} A
foundational model for joint sequence-function multi-species modeling at
scale for long-range genomic prediction. \emph{bioRxiv}
2025.12.22.695963 (2025)
doi:}\href{http://dx.doi.org/10.64898/2025.12.22.695963}{10.64898/2025.12.22.695963}\href{http://paperpile.com/b/FLFbSd/zg9K}{.}

6. \href{http://paperpile.com/b/FLFbSd/5VVM}{Zhou, Z. \emph{et al.}
DNABERT-2: Efficient Foundation Model and Benchmark For Multi-Species
Genome. (2023).}

7. \href{http://paperpile.com/b/FLFbSd/j27Z}{Brixi, G. \emph{et al.}
Genome modelling and design across all domains of life with Evo 2.
\emph{Nature} \textbf{652}, 1349--1361 (2026).}

8. \href{http://paperpile.com/b/FLFbSd/HqTS}{Ma, M. \emph{et al.}
HybriDNA: A Hybrid Transformer-Mamba2 Long-Range DNA Language Model.
(2025).}

9. \href{http://paperpile.com/b/FLFbSd/0vKI}{Shao, B. \& Yan, J. A
long-context language model for deciphering and generating bacteriophage
genomes. \emph{Nature Communications} \textbf{15}, 9392 (2024).}

10. \href{http://paperpile.com/b/FLFbSd/h0C1}{Yu, M., Wang, D., Shan,
Q., Reed, C. J. \& Wan, A. The Super Weight in Large Language Models.
(2024).}

11. \href{http://paperpile.com/b/FLFbSd/XSIF}{The Spike, the Sparse and
the Sink: Anatomy of Massive Activations and Attention Sinks.}
\url{https://arxiv.org/html/2603.05498v1}\href{http://paperpile.com/b/FLFbSd/XSIF}{.}

\end{document}
